% ECONOMICS_PROBLEM: {"id": "L17-432", "title": "Open versus proprietary AI", "jel_code": "L17", "source_id": "OP-0417"}
% !TeX program = xelatex
\documentclass[11pt,a4paper]{article}
\usepackage[margin=23mm]{geometry}
\usepackage{fontspec}
\setmainfont{Latin Modern Roman}
\usepackage{amsmath,amssymb,mathtools}
\usepackage{xcolor,enumitem,booktabs,longtable,array,xurl,hyperref}
\definecolor{muted}{HTML}{53636C}
\hypersetup{colorlinks=true,urlcolor=blue,linkcolor=blue}
\setlength{\parindent}{0pt}
\setlength{\parskip}{5pt}
\newcommand{\R}{\mathbb R}
\newcommand{\E}{\mathbb E}
\newcommand{\N}{\mathbb N}
\newcommand{\ind}{\mathbf 1}
\newcommand{\Var}{\operatorname{Var}}
\newcommand{\Cov}{\operatorname{Cov}}
\newcommand{\supp}{\operatorname{supp}}
\DeclareMathOperator*{\argmin}{arg\,min}
\DeclareMathOperator*{\argmax}{arg\,max}
\newcommand{\entrymeta}[3]{{\sffamily\small\color{muted}\textbf{\texttt{#1}}\quad|\quad #2\par #3\par}}
\newcommand{\Problemref}[1]{\texttt{#1}}
\newcommand{\JELref}[2]{\texttt{#2} (JEL \texttt{#1})}
\begin{document}
\section*{Open versus proprietary AI}
\textbf{Problem ID:} \texttt{L17-432}\quad \textbf{JEL:} \texttt{L17}\par
% BEGIN ECONOMICS STATEMENT
\label{problem:OP-0417}
{\sffamily\small\color{muted}Original subject group: Innovation and AI economics\par}
\label{definitions:problem:30007699}
\textbf{Audit classification:} Background; proposed extension. The 2026 survey metadata and abstract are verified. Full text was unavailable and no exact open-model innovation passage was checked; this causal comparison remains editorial.
\subsection*{Definitions and precise research target}
Consider downstream AI startups in prespecified product markets over three years. At baseline, compare two access regimes for models of matched measured capability: \(O\), which permits downloading, modification and redistribution under a stated license, and \(P\), which provides a priced proprietary service. Compute access and initial support are recorded separately. For market \(m\), let \(N_m(r)\) be independently validated new-product value under regime \(r\in\{O,P\}\), and \(C_m(r)\) all real development, adaptation and deployment costs. Define the downstream innovation effect \(\tau_N=E[N_m(O)-N_m(P)]\) and net-value effect \(\tau_W=E[N_m(O)-C_m(O)-N_m(P)+C_m(P)]\). Expectations concern eligible markets, not all AI development. Clustered rollout, credible licensing changes or another specified design must address endogenous selection into regimes and model-quality differences. Measure knowledge spillovers, entrant survival and incumbent responses, and separately include the upstream developer's future training investment when evaluating a sustained regime. Determine which combinations of appropriability, modification needs and complementary resources make \(\tau_W>0\), and whether short-run downstream gains persist once upstream incentives adjust. Existing theories of open and proprietary models support the comparison; this particular empirical target and general-equilibrium extension are proposed specifications.

\textbf{Definitions, assumptions and mathematical qualifications.} The open regime specifies weights access, modification, redistribution and license restrictions; an API is a different treatment even at equal baseline benchmark capability. Initial capability and compute allowance are preregistered baseline attributes. Capability subsequently produced by the regime is an outcome, so conditioning on it may remove part of the causal effect. \(N_m(r)\) is monetized validated product value including unsuccessful ventures, and \(C_m(r)\) counts resource costs rather than transfers. Define downstream \(\tau_W\) separately from system welfare, which additionally includes upstream training costs and benefits. A sustained regime requires an upstream innovation and model-release law and common equilibrium selection.
\subsection*{Variants and assumption changes}
\begin{enumerate}
\item \textbf{Editorial, fixed upstream technology.} Compare downloadable models and priced services matched at baseline capability, randomizing compute support separately; estimate the downstream three-year effect.
\item \textbf{Editorial, sustained system.} Endogenize future upstream training, release decisions and spillovers under the two license regimes. Compare system welfare rather than append upstream spending mechanically to an unchanged downstream forecast.
\end{enumerate}
\subsection*{References and source audit}
\textbf{Checked source.} Annual Review of Economics 18 (2026), pp. 417--445, abstract; advance publication 6 May 2026. The abstract surveys AI competition and interacting bottlenecks. Publisher metadata and abstract retrieved; dynamic full text unavailable. \url{https://www.annualreviews.org/content/journals/10.1146/annurev-economics-051624-061832}.
\begin{enumerate}\raggedright
\item\hypertarget{definitions:source:30007699:1}{} Joshua S. Gans. 2026. \emph{Market Power in Artificial Intelligence}. Annual Review of Economics 18 (2026), pp. 417--445, abstract; advance publication 6 May 2026.
\url{https://www.annualreviews.org/content/journals/10.1146/annurev-economics-051624-061832}.
DOI: \url{https://doi.org/10.1146/annurev-economics-051624-061832}.
\end{enumerate}

\subsection*{Common conventions: Source mathematical conventions}

These conventions apply to each entry unless explicitly replaced.

\textbf{Models and identification.} A primitive model \(m\) specifies all probability laws, feasible choices, information, equilibrium and measurement rules used by its target. The admissible class \(\mathcal M\) is fixed independently of outcomes. If \(P_O(m)\) is its observable probability law and \(\tau(m)\) the target, its identified set at observable law \(P\) is
\[
 \mathcal I_\tau(P)=\{\tau(m):m\in\mathcal M,\ P_O(m)=P\}.
\]
Point identification means that this set is a singleton. An empty set signals incompatibility of data and assumptions. Coordinate bounds need not describe the joint set. Bounds are sharp only if every retained target is attainable or arbitrarily approachable by compatible models. Finite bounds require explicit restrictions on unobserved quantities; unrestricted confounding, hidden wealth, extrapolation or arbitrary equilibrium selection can yield uninformative sets.

\textbf{Causal interventions.} A potential outcome \(Y_i(p)\) is the outcome under a complete policy regime \(p\), including timing, financing, information and market spillovers. The target population and units are fixed before treatment. With interference, use the complete assignment vector or a justified exposure mapping; individual no-interference notation alone cannot identify an economy-wide intervention. A difference of mean potential outcomes does not require the cross-world joint distribution. Conditioning on post-treatment migration, survival, enrollment or employment changes the estimand and needs separate justification.

\textbf{Statistical statements.} An estimator, confidence set, forecast or test specifies a sampling law, model class and loss/coverage criterion. Uniform \((1-\alpha)\)-coverage means \(\inf_{m\in\mathcal M}\Pr_m\{\tau(m)\in C_n\}\geq1-\alpha\), or an explicitly asymptotic version. The whole parameter space trivially has coverage, so informativeness requires a stated width, risk or power objective. A forecasting improvement is relative to a named benchmark, information set and evaluation population.

\textbf{Equilibrium and optimization.} An equilibrium comprises feasible plans, optimality, beliefs where needed, budget constraints and all market/resource clearing. The primitives or a stated benchmark define each correspondence \(\mathcal E(p;m)\). Nonemptiness is assumed or proved, never inferred just from compact policy sets. A selected-equilibrium target uses a declared selection rule; a robust target quantifies over all permitted equilibria. Use suprema or infima unless attainment is established. An \(\varepsilon\)-optimal policy is feasible and within \(\varepsilon\) of the finite optimum value. Report an empty feasible set separately.

\textbf{Welfare and accounting.} Normative utility, interpersonal normalization, social weights, numeraire, discounting and the use of public receipts are primitives. Behavioral utility need not coincide with the welfare criterion. Household welfare includes ownership income and rents once; adding profits again to compensating variation generally double-counts them. A monetary equivalent is defined by an explicit transfer and price convention, with monotonicity and range conditions. Average effects, mechanism contributions, transfers and welfare losses are different targets. Infinite sums require convergence and dynamic feasible plans require terminal or no-Ponzi conditions.

\textbf{Source versions.} A working-paper date, revision date and journal publication year are distinguished. A DOI for a journal version is not automatically the DOI of an earlier manuscript. Locators refer to the stated version and printed pages where specified. A metadata-only check does not verify a claim about a full-text conclusion. Every entry's source subsection narrows the provenance claim accordingly.
% END ECONOMICS STATEMENT
\end{document}
